\centerline{\bf \Huge Зимняя Сессия}
\centerline{\bf \Huge }

\begin{flushright}
        \scriptsize{\textbf{12:31. Идёт снег. Прибыла Елена Михайловна. Начало Зимней Сессии. Выживут сильнейшие}}
\end{flushright}


\problem{1} Совсем скоро наступит Новый 2025 год. Чему равна сумма всех нечётных годов, начиная с 1 и заканчивая 2025?

\problem{2} Недавно Елена Михайловна собирала вкусные подарки для своих учеников и заметила, что 3/25 от общей массы подарков равны 243 кг. Чему равна общая масса всех подарков?

\problem{3} Анджур Аркадьевич \textit{серьёзный игрок} в бравл старс и, как любой \textit{серьезный игрок}, любит немного преувеличить имеющееся количество кубков. 

Если он вам скажет: «У меня 10000 кубков»,

То на самом деле у него 7000 кубков :)

а) Сколько на самом деле кубков у Анджура Аркадьевича, если он вам сказал «У меня 25000 кубков»?

б) У Анджура Аркадьевича на самом деле 21000 кубков. Какое количество кубков скажет вам Анджур Аркадьевич? 

\problem{4} Найдите значение выражения:

$\Bigr(2025\dfrac{1}{11} - 1\dfrac{3}{11} + 3\dfrac{4}{11} - 2\dfrac{2}{11}\Bigl);$

\problem{5} При каком значении $Z$ число $3\dfrac{9}{17}$ является корнем уравнения:

$X + Z - 2\dfrac{5}{17} = 5\dfrac{4}{17};$
